\subsection{Strong fibres and positive root factors}

Adjacent phase preparations must give the same passing condition on
their overlap. We identify their strongly contracting root fibres on
one original transverse section, then divide along the common root
surface. The resulting positive unit and its inverse have bounded
jets because the section coordinates are regular.

\begin{proposition}[Compatibility of prepared slow factors]
\label{rw:phase-compatibility}
Consider two of the finite-order phase preparations on a common smaller
original tube and a common initial section. Assume their input derivatives
have nonzero margins, with compatible orientations, and their positive
time changes have bounded ratios and inverses. Use the normal domination
chosen in Section~\ref{rw:transitions} for the prescribed finite order.
Then the root fibres of the two preparations agree. Near a double root,
their pulled-back quadratic slow factors differ by a positive unit with
bounded finite jets and inverse jets. The same conclusion holds after
removing the common invariant-axis factor at a weak or horizontal cubic
boundary. The unit estimates use one additional derivative for each
division performed.
\end{proposition}

\begin{proof}
We first identify the root fibre at a positive scale. If $x=r$ is an
equilibrium of the scalar phase equation, its normal fibre consists of
orbits approaching that equilibrium at the normal contraction rate.
Conversely, for an orbit with $x\ne r$, uniqueness for the scalar
equation and its Lipschitz bound give
\[
 |x(t)-r|\ge |x(0)-r|e^{-C\varpi t}.
\]
Here $\varpi=\epsilon$ for the lower phase and is the selected small
angular width for the upper preparation. A regular phase projection is
Lipschitz in the original state. Such an orbit cannot approach the
equilibrium at the faster normal rate. Thus the strongly contracting
orbit set is exactly the root fibre of either preparation. Bounded
positive time changes preserve the strict separation after the widths
are chosen to include their bounds. This proves equality of the two
root fibres in the original tube, within the unextended field. The
finite-smooth parameter limit gives their identity at zero scale.

Choose the first Morse input $s$ and its gap $\mu$ so that its factor
is $s^2+\mu$. Pull the second preparation back to that same section:
\[
 F_2(s,\mu,\mathbf c)=h(s,\mu,\mathbf c)^2+\eta_2(\mu,\mathbf c).
\]
Here $\mathbf c$ denotes the passive coefficients and $\eta_2$ is
the second slow gap, independent of $s$ and distinct from $\nu$.
At the common double root $h=\eta_2=0$ and $|h_s|>c>0$.
Equality of root fibres gives the exact identity
\begin{equation}\label{rw:centers:root-identity}
 F_2(s,-s^2,\mathbf c)=0.
\end{equation}
Its second derivative at the double root is
\[
 (F_2)_{ss}-2(F_2)_\mu=0,\qquad
 (F_2)_{ss}=2h_s^2.
\]
Consequently $(F_2)_\mu=h_s^2>0$ there, with a uniform positive
margin on a smaller coefficient block. Choose a product box
$|s|<r$, $|\mu|<2r^2$ inside that block. The segment between
$-s^2$ and $\mu$ stays in the larger box, and
\begin{equation}\label{rw:centers:root-division}
 F_2=(s^2+\mu)U,\qquad
 U=\int_0^1(F_2)_\mu
       (s,-s^2+\theta(\mu+s^2),\mathbf c)\,\dd\theta>u_*>0.
\end{equation}
Differentiation of this integral gives the finite jets of $U$;
the positive margin gives those of $U^{-1}$. This also proves the
comparison on the positive-gap side, where neither factor has a root.

For a weak factor $u(\gamma+u)$, the common invariant axis and the
regular orientation give a second input $v=u h(u,\gamma,\mathbf c)$,
where $h>h_*>0$. Its full pulled-back factor is
$u h(\gamma_2+u h)$, up to a positive time unit. Remove the factor
$u$ first. The other root gives
$\gamma_2(-u,\mathbf c)+u h(u,-u,\mathbf c)=0$.
At $u=\gamma=0$, differentiation shows
$\partial_\gamma\gamma_2=h(0,0,\mathbf c)$; hence the derivative
in $\gamma$ of the divided factor is $h(0,0,\mathbf c)^2>0$.
Integration from $\gamma=-u$ now gives a positive unit times
$\gamma+u$, as in \eqref{rw:centers:root-division}. A bounded positive
time unit preserves its sign and finite jets.

At the horizontal cubic boundary, first remove the invariant factor
$x$. The first remaining factor is the exact $F_I$ of
\eqref{rw:transitions:lower-factor}, with
$\partial_{\bar B}F_I=1+O(\epsilon)>1/2$ near the origin.
It has a unique root surface
$\bar B=\mathcal B(x,\bar a,\mathbf c,\epsilon)$.
At the cubic point the two divided factors have positive quadratic
coefficients, since the state and time orientations are positive.
At zero scale and $\bar a=0$, the first surface is $\bar B=-x^2$.
Twice differentiating the second divided factor on this surface gives
its positive transverse $\bar B$ derivative, exactly as above.
Continuity preserves that margin for small scale and coefficients.
Divide each factor in $\bar B$ from the same surface $\mathcal B$;
the ratio of their two positive integral units proves the assertion.
All these segments are chosen in one larger coefficient box before
restricting to actual passing points. A simple root uses the same
one-variable division in its nondegenerate splitting coordinate.
\end{proof}

Use the finite-root comparison where the time-change ratio is bounded;
use the horizontal weak/cubic coordinates when a root escapes to
infinity. These comparisons identify the slow passing conditions and
their regular factors. To form a logarithmic derivative ratio, both
paths must also end in the same original central-height coordinate.
